\documentclass[10pt,a4paper]{amsart}
\usepackage[margin=19mm]{geometry}
\usepackage{amsmath,amssymb,mathtools}
\usepackage[scr=rsfs,cal=euler]{mathalfa}
\usepackage{xfrac}
\usepackage[unicode,hidelinks,bookmarks=false]{hyperref}
\newcommand{\ie}{i.e.\ }
\newcommand{\cf}{cf.\ }
\newcommand{\resp}{respectively\ }
\newcommand{\Subsection}{\subsection}
\providecommand{\nobreakdash}{\nobreak\hbox{-}}
\setlength{\emergencystretch}{2em}
\multlinegap0pt
\usepackage{bm}
\usepackage{bbm}
\usepackage{dsfont}
\usepackage{xspace}
\usepackage{cancel}
\usepackage{mathtools}
\usepackage{amsthm}
\makeatletter
\@ifundefined{theorem}{\newtheorem{theorem}{Theorem}}{}
\@ifundefined{definition}{\newtheorem{definition}[theorem]{Definition}}{}
\@ifundefined{lem}{\newtheorem{lem}{Lemma}}{}
\@ifundefined{rem}{\newtheorem{rem}{Remark}}{}
\makeatother
\def\I{{\mathbb {I}}}                       \def\M{{\mathbb {M}}}
\def\PD{{\mathcal P}}                     \def\PDW{{\tilde{\cal P}}}
\def\P{{\mathbb {P}}}                       \def\Q{{\mathbb {Q}}}
\title[Convergence of sequences of ordered selections]{Convergence of sequences of ordered selections}
\author{B. Fazekas, I. Fazekas}
\date{Source: arXiv:2511.12800v1 (CC BY 4.0)}
\begin{document}
\maketitle

\noindent\emph{Source and licence.} Convergence of sequences of ordered selections. Version arXiv:2511.12800v1 (CC BY 4.0), distributed under the Creative Commons Attribution 4.0 International licence, as recorded in the arXiv licence metadata for this exact version and shown on the abstract page https://arxiv.org/abs/2511.12800v1. Full rights notice: licenses/arxiv-2511.12800-CC-BY-4.0.txt.\par\smallskip

\noindent\textit{Editorial scope.} Counted unit: the Lem whose source label is \texttt{H4.2} in the article (source0.tex lines 559--566, source label \texttt{H4.2}), delivered together with the article's own complete original proof of it (source0.tex lines 569--763, one \texttt{proof} environment of 195 lines). No mathematics is written, repaired or re-derived: statement and proof are the article's own text line for line. Required context retained uncounted: def3 (lines 420--449); def4 (lines 471--498); H7 (lines 505--517); H34 (lines 554--556). Every definition, display and label of those lines is retained unchanged. Omitted from this package and not redistributed: 1--419, 450--470, 499--504, 518--553, 557--558, 567--568, 764--1284. Nothing inside a retained excerpt is omitted or altered: every retained body is delivered exactly as the article's lines are written, so no omission receipt of any kind accompanies this package. Citations: the retained excerpts carry no citation command at all, so no bibliography entry of the article is transcribed and no reference is redistributed. Reference adaptation: none was needed and none was made; every reference inside the delivered unit resolves inside it, so no unresolved reference or citation is produced. Printed slips kept literal: no printed word, symbol, hypothesis or number of the counted statement or of its proof was altered. Layout and class: the article's own class is \texttt{article} with packages including amsmath, amssymb, amsthm, multirow, graphicx, color, subcaption, wrapfig, multicol; the wrapper is the lane's pinned \texttt{amsart} editorial wrapper at 10pt on A4, which restates the theorem-like environments the retained text needs and adds the article's own macro definitions that the retained text needs (\texttt{\textbackslash I}, \texttt{\textbackslash P}, \texttt{\textbackslash PD}), with extra wrapper packages (bm, bbm, dsfont, xspace, cancel, mathtools). The delivered numbering is the unit's own. Assets: no retained span contains a figure environment, a graphics-inclusion command, a PDF-page inclusion, a table or a TikZ picture; no image file is copied into this package, the package declares no asset dependency and no image file is redistributed with it. Licence: version arXiv:2511.12800v1 of this article is distributed under the Creative Commons Attribution 4.0 International licence, as recorded in the arXiv licence metadata for this exact version and shown on the abstract page https://arxiv.org/abs/2511.12800v1. A full rights notice is delivered at licenses/arxiv-2511.12800-CC-BY-4.0.txt. Characters: every retained excerpt is pure ASCII, so the delivered engine, XeLaTeX with its default Latin Modern text font, reports no missing character.
\begin{definition}  \label{def3}
Let $\mu$ be a $\lambda$-permuton with marginal distribution functions 
$F_x$, $F_y$ and parameter $\lambda$.
Let $k$ be a fixed positive integer.
We define the $k$-subdivision of $\mu$ as the following array of rectangles $R_{i,j}$.
Let
$$
0=x_0 < x_1 < \dots < x_k=\lambda, \qquad
0=y_0 < y_1 < \dots < y_k=1
$$
so that $x_i$ is the $\frac{i}{k}$-quantile of $F_x$ ($i=1,\dots ,k-1$)
and
$y_j$ is the $\frac{j}{k}$-quantile of $F_y$ ($j=1,\dots ,k-1$).
Then define the rectangle $R_{i,j}$ as
$$
R_{i,j} = [x_{i-1}, x_i ) \times [y_{j-1}, y_j ), \quad i,j= 1, 2, \dots , k.
$$
Let $\sigma=(\sigma(1), \sigma(2), \dots , \sigma(k))\in \PD_k$ be a permutation of $[k]$.
Define the two-dimensional density function $f_\sigma$ as
\begin{equation}
f_\sigma (x,y)= \frac{1}{k} \cdot \frac{1}{x_i- x_{i-1}} \cdot \frac{1}{y_{\sigma(i)}- y_{\sigma(i)-1}},
\ \ \text{if} \ \ (x,y) \in R_{i,\sigma(i)}
\end{equation}
for some $i= 1, \dots , k$
and let $f_\sigma (x,y)=0$ otherwise.
The corresponding distribution function is denoted by $ F_\sigma (x,y)$.
%
The density function $f_\sigma$ defines a probability measure $\mu_\sigma$ on  
$[0, \lambda] \times [0,1]$.
\end{definition}

\begin{definition}  \label{def4}
Let $\mu$ be a generalized permuton.
Choose a sample from the distribution $\mu$, i.e.
let $(\xi_1, \eta_1), \dots , (\xi_k, \eta_k)$
be independent identically distributed random vectors such that the distribution
of $(\xi_i, \eta_i)$ is $\mu$ for each $i\in [k]$.
Let
$$
\xi_1^\ast< \xi_2^\ast < \dots < \xi_k^\ast , \qquad
\eta_1^\ast< \eta_2^\ast < \dots < \eta_k^\ast
$$
be the ordered samples
(remark that the marginal distribution functions are continuous, so the sample elements are
distinct with $\mu$-probability $1$, so we can use strict inequalities in the above relations).
%
The  $\mu$-random permutation is defined as
$$
\sigma=(\sigma(1), \sigma(2), \dots , \sigma(k)),
$$
if for any $i\in[k]$ we have $\xi_i^\ast =\xi_l$ and $\eta_l = \eta_{\sigma(i)}^\ast$ for 
a certain $l$.
%
We shall denote the above $\mu$-random permutation by $\mu^{(k)}$.
%
Now, the $\mu$-random generalized permuton is defined as
$\sigma(k,\mu) =\mu_\sigma$, where $\sigma$ is a $\mu$-random permutation 
and $\mu_\sigma$ is from Definition \ref{def3}.
\end{definition} 

\begin{rem}
Let $F$ be a two-dimensional distribution function with marginal distribution functions
$F_x$  and $F_y$.
Then for $x_1<x_2$ and $y_1<y_2$ we have
\begin{equation} \label{H7}
F(x_2, y_2) -  F(x_1, y_1) \le (F_x(x_2) -  F_x(x_1)) + (F_y(y_2) -  F_y(y_1)).
\end{equation}
If the slopes of $F_x$  and $F_y$ are bounded by $C$, then
for $x_1<x_2$ and $y_1<y_2$ we have
\begin{equation} \label{H7+}
F(x_2, y_2) -  F(x_1, y_1) \le C((x_2 -  x_1) + (y_2 -  y_1)).
\end{equation}
\end{rem}

\begin{equation} \label{H34}
d_\infty(\mu_1,\mu_2) \le d_\square (\mu_1,\mu_2) \le 4 d_\infty(\mu_1,\mu_2) .
\end{equation}

\begin{lem}  \label{H4.2}
There exists a $k_0$ such that for $k> k_0$ and for any $\lambda$-permuton $\mu$ we have
\begin{equation} \label{H37}
\P \left( d_\square \left(\mu, \sigma(k, \mu)\right) \le 16 k^{-\frac{1}{4}}   \right)
\ge 1- \frac{1}{2} e^{-\sqrt{k}},
\end{equation}
where $\sigma(k, \mu)$ is an arbitrary $\mu$-random $\lambda$-permuton introduced in Definition \ref{def4}.
\end{lem}

\begin{proof}
Let $\mu$ be a $\lambda$-permuton on $[0,\lambda] \times [0,1]$ and let 
$F$ be its two-dimensional distribution function.
Let $\sigma(k, \mu)$ be a $\mu$-random $\lambda$-permuton and let $F_k$ be its 
two-dimensional distribution function.
By \eqref{H34}, if we prove
\begin{equation} \label{H39}
\P \left( \sup_{x,y\in[0,1]} |F(x,y) -F_k(x,y)| \ge 4 k^{-\frac{1}{4}}   \right)
\le \frac{1}{2} e^{-\sqrt{k}},
\end{equation}
then we obtain \eqref{H37}.

Let us use the notation from Definition \ref{def3}.
Let $(x,y)\in R_{i,j}$.
Then
\begin{eqnarray} \label{H18}
&&| F_k(x,y)- F(x,y) |\le \\
&\le & | F_k(x,y)- F_k(x_i,y_j) |+ | F_k(x_i,y_j)- F(x_i,y_j) | +| F(x_i,y_j)- F(x,y) |\le  \nonumber \\
&\le & \frac{2}{k}+ | F_k(x_i,y_j)- F(x_i,y_j) | + \frac{2}{k}, \nonumber
\end{eqnarray}
using \eqref{H7} and the fact that $x_i$ is the $\frac{i}{k}$ quantile of $F_x$ and $y_j$ is the $\frac{j}{k}$ quantile of $F_y$.

Let $\varepsilon= k^{-\frac{1}{4}}$.
Let $(\xi_1, \eta_1), \dots , (\xi_k, \eta_k)$ be the sample from the distribution $\mu$.
Then we have
\begin{equation} \label{H41-}
F_k(x_i,y_j) = \frac{1}{k} \sum_{t=1}^k \I\{ \xi_t \le \xi_i^\ast, \eta_t \le \eta_j^\ast \} ,
\end{equation} 
where $\I$ denotes the indicator function.
%%%%%%%%%%%%%%%%%%%%%%%%%%%%%%%%%%%%%%%%%%%%%%%%%%%%%%%%%%%%%%%%%%%%%%%%%%%%%%%%5
Now, we shall show that for all $i,j\in [k]$
\begin{equation} \label{H41}
\P \left( F_k(x_i,y_j) > F(x_i,y_j) +3\varepsilon   \right)
< 3 e^{-2\sqrt{k}} .
\end{equation} 
Let
$x'_i$ and $y'_j$ be such that
\begin{equation*}
F_x(x_i+ x'_i) = \frac{i}{k }  + \varepsilon, \qquad  F_y(y_j+ y'_j) = \frac{j}{k }  + \varepsilon .
\end{equation*} 
So \eqref{H7} gives
\begin{equation} \label{H42}
F(x_i+ x'_i, y_j+y'_j) +\varepsilon \le F(x_i, y_j) + 3 \varepsilon .
\end{equation} 
Let $A$ be the following event
\begin{equation} \label{H42+}
A=  \left\{\frac{1}{k} \sum_{t=1}^k \I\{ \xi_t \le \xi_i^\ast, \eta_t \le \eta_j^\ast \} 
> F(x_i+ x'_i, y_j+y'_j) +\varepsilon  \right\} .
\end{equation} 
Then \eqref{H42} gives that the left-hand side of \eqref{H41} is not greater than $\P(A)$.
So we have to find an upper bound of $\P(A)$.
We have
\begin{eqnarray} \label{H43}
\P(A) &=& \P\left(A \cap \{ \xi^\ast_i <  x_i+ x'_i  \} \cap \{ \eta^\ast_j <  y_j+ y'_j  \} \right) \\
&+&\P\left(A \cap \left(\{ \xi^\ast_i \ge  x_i+ x'_i  \} 
\cup \{ \eta^\ast_j \ge  y_j+ y'_j  \} \right) \right) \nonumber \\
&\le&
\P\left(A \cap \{ \xi^\ast_i <  x_i+ x'_i  \} \cap \{ \eta^\ast_j <  y_j+ y'_j  \} \right) \nonumber \\
&+&\P\left(\xi^\ast_i \ge  x_i+ x'_i  \right) +
\P\left( \eta^\ast_j \ge  y_j+ y'_j  \} \right) . \nonumber 
\end{eqnarray}

We need the following large deviation result for the binomial distribution.
Let $\zeta$ be a binomial random variable with parameters $k$ and $p$.
Then for $\varepsilon>0$
\begin{equation} \label{H44}
\P\left(\frac{1}{k} \zeta \ge p+\varepsilon \right)
\le \exp (-2k \varepsilon^2 ), \ \
\P\left(\frac{1}{k} \zeta \le p-\varepsilon \right)
\le \exp (-2k \varepsilon^2 ) .
\end{equation}
Then, with $p= F_x(x_i+ x'_i) = \frac{i}{k }  + \varepsilon$, we have
\begin{eqnarray*}
&&\P\left(\xi^\ast_i \ge  x_i+ x'_i  \right) \\
&&=\P \left( \frac{1}{k} \sum_{t=1}^k \I\{ \xi_t \le x_i+ x'_i\}  \le \frac{i}{k} \right)
= \P\left(\frac{1}{k} \zeta \le p-\varepsilon \right)
\le \exp (-2k \varepsilon^2 ) .
\end{eqnarray*}
Similarly 
$$
\P\left(\eta^\ast_i \ge  y_j+ y'_j  \right) 
\le \exp (-2k \varepsilon^2 ) .
$$
%
Moreover,
\begin{eqnarray*} 
&& \P\left(A \cap \{ \xi^\ast_i <  x_i+ x'_i \} \cap \{ \eta^\ast_j <  y_j+ y'_j \} \right)  \\
 = &&
\P \Bigg( \left\{ \frac{1}{k} \sum\nolimits_{t=1}^k \I\{ \xi_t \le \xi_i^\ast, \eta_t \le \eta_j^\ast \} 
> F(x_i+ x'_i, y_j+y'_j) +\varepsilon  \right\}     \\
&& \cap \{ \xi^\ast_i <  x_i+ x'_i  \} \cap \{ \eta^\ast_j <  y_j+ y'_j  \} \Bigg)  \\
=&&
\P \Bigg( \left\{ \frac{1}{k} \sum\nolimits_{t=1}^k 
\I\{ \xi_t \le x_i+ x'_i, \eta_t \le y_j+ y'_j \} 
\ge  F(x_i+ x'_i, y_j+y'_j) +\varepsilon  \right\}  \Bigg)   \\
= &&  \P\left(\frac{1}{k} \zeta \ge p+\varepsilon \right)
\le \exp (-2k \varepsilon^2 ) .
\end{eqnarray*}
From the above inequalities
$$
\P(A) \le 3 \exp (-2k \varepsilon^2 ) = 3 \exp (-2k (k^{-\frac{1}{4}})^2 ) = 3 \exp (-2\sqrt{k}).
$$
%
%%%%%%%%%%%%%%%%%%%%%%%%%%%%%%%%%%%%%%%%%%%%%%%%%%%%%%%%%%%%%%%%%%%%%%%%%%
Now, we have to prove that for all $i,j\in [k]$
\begin{equation} \label{H41+}
\P \left( F_k(x_i,y_j) < F(x_i,y_j) -3\varepsilon   \right)
< 3 e^{-2\sqrt{k}} .
\end{equation} 
Let now
$x''_i$ and $''_j$ be such that
\begin{equation*}
F_x(x_i- x''_i) = \frac{i}{k }  - \varepsilon, \qquad  F_y(y_j- y''_j) = \frac{j}{k }  - \varepsilon .
\end{equation*} 
So \eqref{H7} gives
\begin{equation} \label{H42++}
F(x_i- x''_i, y_j-y''_j) -\varepsilon \ge F(x_i, y_j) - 3 \varepsilon .
\end{equation} 
Let $A''$ be the following event
\begin{equation} \label{H42+++}
A''=  \left\{\frac{1}{k} \sum_{t=1}^k \I\{ \xi_t \le \xi_i^\ast, \eta_t \le \eta_j^\ast \} 
< F(x_i- x''_i, y_j-y''_j) -\varepsilon  \right\} .
\end{equation} 
Then \eqref{H42++} gives that the left-hand side of \eqref{H41+} is not greater than $\P(A'')$.
So we have to find an upper bound of $\P(A'')$.
We have
\begin{eqnarray} \label{H43+}
\P(A'') &=& \P\left(A'' \cap \{ \xi^\ast_i >  x_i- x''_i  \} \cap 
\{ \eta^\ast_j >  y_j- y''_j  \} \right) \\
&+&\P\left(A'' \cap \left(\{ \xi^\ast_i \le  x_i- x''_i  \} 
\cup \{ \eta^\ast_j \le  y_j- y''_j  \} \right) \right) \nonumber \\
&\le&
\P\left(A \cap \{ \xi^\ast_i >  x_i- x''_i  \} \cap \{ \eta^\ast_j >  y_j- y''_j  \} \right) \nonumber \\
&+&\P\left(\xi^\ast_i \le  x_i- x''_i  \right) +
\P\left( \eta^\ast_j \le  y_j- y''_j  \} \right) . \nonumber 
\end{eqnarray}
%
We again use the large deviation result for the binomial distribution.
Then
\begin{eqnarray*}
&&\P\left(\xi^\ast_i \le  x_i- x''_i  \right) \\
&&=\P \left( \frac{1}{k} \sum_{t=1}^k \I\{ \xi_t \le x_i- x''_i\}  \ge \frac{i}{k} \right)
= \P\left(\frac{1}{k} \zeta \ge p+\varepsilon \right)
\le \exp (-2k \varepsilon^2 ) .
\end{eqnarray*}
Similarly 
$$
\P\left(\eta^\ast_i \le  y_j- y''_j  \right) 
\le \exp (-2k \varepsilon^2 ) .
$$
Then
\begin{eqnarray*} 
&& \P\left(A'' \cap \{ \xi^\ast_i >  x_i- x''_i \} \cap \{ \eta^\ast_j >  y_j- y''_j \} \right)  \\
 = &&
\P \Bigg( \left\{ \frac{1}{k} \sum\nolimits_{t=1}^k \I\{ \xi_t \le \xi_i^\ast, \eta_t \le \eta_j^\ast \} 
< F(x_i- x''_i, y_j-y''_j) -\varepsilon  \right\}     \\
&& \cap \{ \xi^\ast_i >  x_i- x''_i  \} \cap \{ \eta^\ast_j >  y_j- y''_j  \} \Bigg)  \\
=&&
\P \Bigg( \left\{ \frac{1}{k} \sum\nolimits_{t=1}^k 
\I\{ \xi_t \le x_i- x''_i, \eta_t \le y_j- y''_j \} 
\le  F(x_i- x''_i, y_j-y''_j) -\varepsilon  \right\}     \\
= &&   \P\left(\frac{1}{k} \zeta \le p-\varepsilon \right)
\le \exp (-2k \varepsilon^2 ) .
\end{eqnarray*}
From the above inequalities
$$
\P(A'') \le 3 \exp (-2k \varepsilon^2 ) = 3 \exp (-2\sqrt{k}).
$$
By \eqref{H41} and \eqref{H41+}, we have
\begin{equation} \label{H40}
\P \left( |F_k(x_i,y_j) - F(x_i,y_j)| > 3\varepsilon   \right)
\le 6 e^{-2\sqrt{k}} .
\end{equation} 
Let $k$ be so large that $\frac{4}{k} <k^{-\frac{1}{4}}$.
Then, by \eqref{H18},
$$
| F_k(x,y)- F(x,y) |\le   k^{-\frac{1}{4}}+ | F_k(x_i,y_j)- F(x_i,y_j) | 
$$
if $(x,y) \in R_{i,j}$.
Then, for large enough $k$
\begin{eqnarray*} 
&&\P \left( \sup_{x,y\in[0,1]} |F(x,y) -F_k(x,y)| \ge 4 k^{-\frac{1}{4}}   \right) \\
&\le & 
\P \left( \max_{i,j\in[k]} |F(x_i,y_j) -F_k(x_i,y_j)| > 3 k^{-\frac{1}{4}}   \right) \\
&\le & 
\sum_{i,j\in[k]}
\P \left( |F(x_i,y_j) -F_k(x_i,y_j)| > 3 k^{-\frac{1}{4}}   \right) \\
&\le & 
6 k^2 e^{-2\sqrt{k}}
= \left( 12 k^2 e^{-\sqrt{k}} \right) \left(\frac{1}{2} e^{-\sqrt{k}}\right) \le
\frac{1}{2} e^{-\sqrt{k}}.     
\end{eqnarray*}
%%%%%%%%%%%%%%%%%%%%%%%%%%%%%%%%%%%%%%%%%%%%%%%%%%%%%%%%%%%%%%%%%%%%%%%%%%
This completes the proof.
\end{proof}

\end{document}
